%----------------------------------------------------------------------------------------
% STELSELS. Alle opgaven zijn eigen werk.
%
% Een hoofdstuk kernvaardigheden (ontwerpregels, deel Kernvaardigheden): beknopte theorie, daarna drill.
% Labels van dit hoofdstuk beginnen met h7-. Notatie zoals in Stelsels (2de graad): \stelsel, koppels (x, y).
%----------------------------------------------------------------------------------------

\mychapter{7}{Stelsels}

\lkr{Herhaling van Stelsels (3de jaar, LPD 26-27). Nodig in Analytische meetkunde (snijpunt van twee rechten, een rechte door twee punten), Vectoren (een vector ontbinden) en in de wetenschapsvakken. Bundel 6 (drie onbekenden) staat in Stelsels alleen als opdracht voor pakket 4-5.}

\section{Wat je moet kennen}

\begin{theorie}
\begin{minipage}[t]{.47\linewidth}
\vspace{0pt}
Een \textbf{\textcolor{accent}{oplossing}} van een stelsel met twee onbekenden is een koppel $(x, y)$ dat aan \emph{beide} vergelijkingen voldoet.

\vspace{1mm}
\textbf{\textcolor{accent}{Substitutie:}} zonder een onbekende af en vul ze in de andere vergelijking in.

\vspace{1mm}
\textbf{\textcolor{accent}{Combinatie:}} maak de coëfficiënten van een onbekende tegengesteld en tel de vergelijkingen op.
\end{minipage}\hfill
\begin{minipage}[t]{.48\linewidth}
\vspace{0pt}
Grafisch is elke vergelijking een rechte:

\vspace{1mm}
snijdende rechten: \emph{één} oplossing;

evenwijdige rechten: \emph{geen} oplossing, het stelsel is \textbf{\textcolor{accent}{strijdig}}, $V = \emptyset$;

samenvallende rechten: \emph{oneindig veel} oplossingen, het stelsel is \textbf{\textcolor{accent}{onbepaald}}.
\end{minipage}
\end{theorie}

\textbf{Klassieke fouten.} Een minteken vergeten als je een vergelijking aftrekt: trek \emph{elke} term af, ook het rechterlid. Alleen $x$ berekenen en $y$ vergeten. Na het oplossen niet controleren: vul het koppel in \emph{beide} vergelijkingen in.

\voorbeeld{}{Los op: \quad $\stelsel{3x + 4y &= 2}{5x - 2y &= 12}$\label{vb:h7-stelsel}}{%
\begin{stappen}
3x + 4y + 10x - 4y = 2 + 24 & Tweede vergelijking maal $2$, dan optellen: $y$ valt weg. \\
13x = 26, \text{ dus } x = 2 & \\
6 + 4y = 2, \text{ dus } y = -1 & Invullen in de eerste vergelijking. \\
(2, -1) & Controle in de tweede: $10 + 2 = 12$.
\end{stappen}}

%========================================================================================
\opdrachtenkop

\begin{oefeningen}
\needspace{10\baselineskip}
\oef{h7-o1}

\needspace{10\baselineskip}
\oef{h7-o2}

\needspace{10\baselineskip}
\oef{h7-o3}

\needspace{10\baselineskip}
\oef{h7-o4}

\needspace{10\baselineskip}
\oef{h7-o5}

\needspace{10\baselineskip}
\oef{h7-o6}

\needspace{10\baselineskip}
\oef{h7-o7}
\end{oefeningen}

%----------------------------------------------------------------------------------------
\needspace{14\baselineskip}
\opdrachtsectie{Test jezelf}

\begin{oefeningen}
\oef{h7-t1}

\oef{h7-t2}

\oef{h7-t3}

\oef{h7-t4}
\end{oefeningen}

\herhaalschema{%
  \fv{oef:h7-t1}{vb:h7-stelsel}{oef:h7-o2}%
  \fzv{oef:h7-t2}{oef:h7-o3}%
  \fzv{oef:h7-t3}{oef:h7-o4}%
  \fzv{oef:h7-t4}{oef:h7-o5}}

\ifzelfstudie
\section*{Maak je eigen samenvatting}
Sluit dit boek en schrijf in je schrift: wat een oplossing van een stelsel is, de methodes substitutie en combinatie, wanneer een stelsel strijdig of onbepaald is, en hoe je controleert. Vergelijk daarna met het kader.
\fi
